% Part: computability
% Chapter: computability-theory
% Section: ce-closed-cup-cap

\documentclass[../../../include/open-logic-section]{subfiles}

\begin{document}

\olfileid[hi]{cmp}{thy}{clo}

\olsection[C.E. समुच्चयों का संघ और प्रतिच्छेद]{संगणनीयतः परिगणनीय
  समुच्चय संघ और प्रतिच्छेद के अंतर्गत संवृत हैं}

अगला प्रमेय संगणनीयतः परिगणनीय समुच्चयों के कुछ संवृतता-गुण देता है।

\begin{thm}
मान लें कि $A$ और $B$ संगणनीयतः परिगणनीय हैं। तब $A \cap B$ और
$A \cup B$ भी ऐसे ही हैं।
\end{thm}

\begin{proof}
\olref[eqc]{thm:ce-equiv} हमें संगणनीयतः परिगणनीय समुच्चयों के विभिन्न
अभिलक्षणों का उपयोग करने देता है। स्पष्टीकरण के लिए हम कुछ भिन्न प्रमाण देंगे।

पहले प्रमाण के लिए मान लें कि $A$ का परिगणन संगणनीय फलन~$f$ और $B$ का
परिगणन संगणनीय फलन~$g$ करता है। मानें
\begin{align*}
h(x) & = \umin{y}{(f(y) = x \lor g(y) = x)} \text{ तथा}\\
j(x) & = \umin{y}{(f((y)_0) = x \land g((y)_1) = x)}.
\end{align*}
तब $A \cup B$, $h$ का प्रांत है और $A \cap B$,~$j$ का प्रांत है।

\begin{explain}
संगणनात्मक पदों में यह हो रहा है: $A$ और $B$ का परिगणन करने वाली
प्रक्रियाएँ मिलने पर हम अर्ध-निर्णय कर सकते हैं कि अवयव $x$, $A \cup B$
में है या नहीं—इसके लिए किसी भी परिगणन में $x$ खोजें; और हम अर्ध-निर्णय
कर सकते हैं कि अवयव $x$, $A \cap B$ में है या नहीं—इसके लिए दोनों
परिगणनों में एक साथ $x$ खोजें।
\end{explain}

दूसरे प्रमाण के लिए फिर मान लें कि $A$ का परिगणन~$f$ और $B$ का
परिगणन~$g$ करता है। मानें
\[
k(x) = \begin{cases}
f(x/2) & \text{यदि $x$ सम है} \\
g((x-1)/2) & \text{यदि $x$ विषम है।}
\end{cases}
\]
तब $k$, $A \cup B$ का परिगणन करता है; विचार यह है कि $k$, $f$ और~$g$
द्वारा दिए परिगणनों को बारी-बारी से लेता है। $A \cap
B$ का परिगणन अधिक
कठिन है। यदि $A \cap B$ रिक्त है, तो वह तुच्छतः संगणनीयतः परिगणनीय है।
अन्यथा $c$ को $A \cap B$ का कोई अवयव मानें और $l$ को परिभाषित करें:
\[
l(x) = \begin{cases}
f((x)_0) & \text{यदि $f((x)_0) = g((x)_1)$} \\
c & \text{अन्यथा।}
\end{cases}
\]
संगणनात्मक पदों में, $l$, $f$ और $g$ के परिगणनों के अवयव-युग्मों से
होकर चलता है और मिलने वाला प्रत्येक मेल निर्गत करता है; अन्यथा वह $c$
निर्गत करके वहीं बना रहता है।

अंतिम प्रमाण के लिए मान लें कि $A$ आंशिक फलन $m(x)$ का \emph{प्रांत}
और $B$ आंशिक फलन $n(x)$ का प्रांत है। तब $A \cap B$ आंशिक फलन
$m(x) +
n(x)$ का प्रांत है।

\begin{explain}
संगणनात्मक पदों में, यदि $A$ उन मानों का समुच्चय है जिन पर $m$ रुकता है
और $B$ उन मानों का समुच्चय है जिन पर $n$ रुकता है, तो $A \cap B$ उन
मानों का समुच्चय है जिन पर दोनों प्रक्रियाएँ रुकती हैं।
\end{explain}

$A \cup B$ को रुकने वाले मानों के समुच्चय के रूप में व्यक्त करना अधिक
कठिन है, क्योंकि $m$ और $n$ का समानांतर अनुकरण करना पड़ता है। $d$ को
$m$ का और $e$ को $n$ का सूचकांक मानें; दूसरे शब्दों में, $m =
\cfind{d}$
और $n = \cfind{e}$। तब $A \cup B$ इस फलन का प्रांत है:
\[
p(x) = \umin{y}{(T(d,x,y) \lor T(e,x,y))}.
\]
\begin{explain}
संगणनात्मक पदों में, आगत $x$ पर $p$, $m$ की रुकने वाली संगणना अथवा $n$
की रुकने वाली संगणना खोजता है और इनमें से कोई मिलते ही रुक जाता है।
\end{explain}
\end{proof}

\end{document}
